\section{附录：场景启用门槛的保守估计表}\label{app:activation-threshold}

本附录给出场景特异可靠度 $r_u^{(s)}$ 的启用门槛 $N_{u,s,\min}$ 在两种备选取值下的保守启用率估计表。门槛最终取值在 PreScreen 结束后冻结。

设核心场景集合为 $\mathcal{S}_{\mathrm{core}}$，其由 \texttt{Calibration\_anchor}+\texttt{Calibration\_core} 第一轮完成后按频率与一致性冻结得到。每位通过 PreScreen 的 worker 在 \texttt{Calibration\_manual} 至少获得 $N_{cal,\min}=25$ 个任务。对每个核心场景 $s$ 统计其频率并取保守下界 $\underline{p}_s$，用二项近似给出启用率下界：
\[\Pr\bigl(N_{u,s}\ge N_{u,s,\min}\bigr)\ge \Pr\bigl(\mathrm{Bin}(N_{cal,\min},\underline{p}_s)\ge N_{u,s,\min}\bigr).\]

\begin{table}[H]
\centering
\caption{启用门槛为 $N_{u,s,\min}=5$ 时的保守启用率估计表}
\label{tab:activation-nusmin-5}
\small
\begin{tabular}{p{0.18\linewidth} p{0.13\linewidth} p{0.13\linewidth} p{0.16\linewidth} p{0.18\linewidth}}
\toprule
\textbf{核心场景 $s$} & $\underline{p}_s$ & $N_{cal,\min}$ & $\mathbb{E}[N_{u,s}]$ & $\Pr\left(N_{u,s}\ge 5\right)$ \\
\midrule
$s_1$ &  & 25 &  &  \\
$s_2$ &  & 25 &  &  \\
$s_3$ &  & 25 &  &  \\
$s_4$ &  & 25 &  &  \\
\bottomrule
\end{tabular}
\end{table}
\FloatBarrier

\begin{table}[H]
\centering
\caption{启用门槛为 $N_{u,s,\min}=6$ 时的保守启用率估计表}
\label{tab:activation-nusmin-6}
\small
\begin{tabular}{p{0.18\linewidth} p{0.13\linewidth} p{0.13\linewidth} p{0.16\linewidth} p{0.18\linewidth}}
\toprule
\textbf{核心场景 $s$} & $\underline{p}_s$ & $N_{cal,\min}$ & $\mathbb{E}[N_{u,s}]$ & $\Pr\left(N_{u,s}\ge 6\right)$ \\
\midrule
$s_1$ &  & 25 &  &  \\
$s_2$ &  & 25 &  &  \\
$s_3$ &  & 25 &  &  \\
$s_4$ &  & 25 &  &  \\
\bottomrule
\end{tabular}
\end{table}

在主报告中同时给出两张表的数值，并披露实际启用比例与退化比例，用于审计场景分层是否在运行中退化为全局。

若第一轮校准后核心场景覆盖不足，则按第~\ref{sec:method-global-reliability}节的覆盖补齐规则，从 \texttt{Calibration\_reserve} 中对核心场景进行补派并再次报告启用率变化。若补齐后总体启用率仍低于 $\tau_{act}$，则主结论采用降级口径，不将场景路由作为强主张。